\documentclass[11pt]{article}

\usepackage[margin=1in]{geometry}
\usepackage{amsmath,amssymb,bm}
\usepackage{siunitx}
\usepackage{hyperref}

\title{Peridynamic Bar With Interface Coupling and Linear Softening Damage}
\author{}
\date{}

\begin{document}

\maketitle

\section{Geometry and Discretization}

The model is a one-dimensional peridynamic bar of length
\begin{equation}
L = 1 .
\end{equation}
The bar is discretized into \(N\) material points. The nodal coordinates are
\begin{equation}
x_i = i\Delta x,
\qquad i = 0,\ldots,N-1,
\end{equation}
where
\begin{equation}
\Delta x = \frac{L}{N-1}.
\end{equation}
Thus the bar occupies
\begin{equation}
x \in [0,L].
\end{equation}

Each point has volume
\begin{equation}
V_i = A\Delta x,
\end{equation}
where the cross-sectional area is currently
\begin{equation}
A = 1.
\end{equation}

The peridynamic horizon is
\begin{equation}
\delta = m\Delta x,
\end{equation}
with
\begin{equation}
m = 3.
\end{equation}
A point \(j\) belongs to the family of point \(i\), denoted \(\mathcal{H}_i\), if
\begin{equation}
0 < |x_j-x_i| \le \delta + \varepsilon_{\mathrm{tol}},
\end{equation}
where a small numerical tolerance is used:
\begin{equation}
\varepsilon_{\mathrm{tol}} = 10^{-12}.
\end{equation}
This tolerance avoids losing horizon bonds due to floating-point roundoff.

\section{Kinematics}

The displacement of material point \(i\) is denoted by
\begin{equation}
u_i .
\end{equation}
The current position is
\begin{equation}
y_i = x_i + u_i.
\end{equation}

For a bond \((i,j)\), the initial bond vector is
\begin{equation}
\xi_{ij} = x_j - x_i,
\end{equation}
with initial length
\begin{equation}
|\xi_{ij}| = |x_j-x_i|.
\end{equation}
The deformed bond vector is
\begin{equation}
\eta_{ij} = y_j-y_i
          = (x_j+u_j)-(x_i+u_i),
\end{equation}
with deformed length
\begin{equation}
|\eta_{ij}| = |y_j-y_i|.
\end{equation}

The bond stretch is
\begin{equation}
s_{ij}
=
\frac{|\eta_{ij}|-|\xi_{ij}|}{|\xi_{ij}|}.
\end{equation}
The one-dimensional bond direction is
\begin{equation}
n_{ij}
=
\frac{\eta_{ij}}{|\eta_{ij}|}.
\end{equation}

For damage, the maximum historical stretch is stored:
\begin{equation}
s^{\max}_{ij}(t)
=
\max_{\tau \le t} s_{ij}(\tau).
\end{equation}

\section{Micromodulus}

The continuum one-dimensional bond-based micromodulus is
\begin{equation}
c_0 = \frac{2E}{A\delta^2}.
\end{equation}
Because the horizon contains a finite number of points, a discrete correction is applied:
\begin{equation}
c = c_0 \frac{m}{m+1}.
\end{equation}
Therefore,
\begin{equation}
c
=
\frac{2E}{A\delta^2}
\frac{m}{m+1}.
\end{equation}
The current nondimensional material parameters are
\begin{equation}
E = 1,
\qquad
A = 1.
\end{equation}

\section{Surface Correction}

Near boundaries, a material point has an incomplete horizon. A geometric correction factor is assigned to each point:
\begin{equation}
G_i
=
\frac{2\delta}{\delta_i^-+\delta_i^+},
\end{equation}
where
\begin{equation}
\delta_i^- = \min(\delta,x_i),
\end{equation}
and
\begin{equation}
\delta_i^+ = \min(\delta,L-x_i).
\end{equation}
Interior points have
\begin{equation}
G_i = 1,
\end{equation}
whereas boundary-near points have
\begin{equation}
G_i > 1.
\end{equation}

For the displacement-prescribed grip nodes, however, the correction is set to unity:
\begin{equation}
G_i = 1.
\end{equation}
This is done for the first and last three nodes. These nodes are externally constrained, so they are not strengthened by the free-surface correction.

For a bond \((i,j)\), the bond correction is
\begin{equation}
G_{ij}
=
\frac{G_i+G_j}{2}.
\end{equation}
The effective integration weight is
\begin{equation}
w_{ij}
=
V_jG_{ij}.
\end{equation}

\section{Elastic Bond Force}

For an intact elastic bond, the force density contribution from point \(j\) to point \(i\) is
\begin{equation}
f_{ij}
=
c\,s_{ij}\,n_{ij}\,w_{ij}.
\end{equation}
The total bar-bar internal force density is
\begin{equation}
p_i^{\mathrm{bar}}
=
\sum_{j\in\mathcal{H}_i}
c\,s_{ij}\,n_{ij}\,w_{ij}.
\end{equation}

\section{Substrate and Interface Coupling}

The substrate is represented by a virtual layer below the bar. The substrate has the same horizontal grid as the bar and is vertically offset by
\begin{equation}
h = \Delta x.
\end{equation}
The substrate displacement is prescribed affinely:
\begin{equation}
u_j^{\mathrm{sub}}
=
\varepsilon x_j.
\end{equation}

For the interaction between bar point \(i\) and virtual substrate point \(j\), the initial interface distance is
\begin{equation}
r_{ij}^{0,\mathrm{sub}}
=
\sqrt{(x_j-x_i)^2+h^2}.
\end{equation}
The current horizontal separation is
\begin{equation}
d_{ij}^{\mathrm{sub}}
=
(x_j+u_j^{\mathrm{sub}})
-
(x_i+u_i).
\end{equation}
The current interface distance is
\begin{equation}
r_{ij}^{\mathrm{sub}}
=
\sqrt{
\left(d_{ij}^{\mathrm{sub}}\right)^2+h^2
}.
\end{equation}
The interface stretch is
\begin{equation}
s_{ij}^{\mathrm{sub}}
=
\frac{
r_{ij}^{\mathrm{sub}}-r_{ij}^{0,\mathrm{sub}}
}{
r_{ij}^{0,\mathrm{sub}}
}.
\end{equation}
The horizontal direction cosine is
\begin{equation}
n_{ij}^{\mathrm{sub}}
=
\frac{
d_{ij}^{\mathrm{sub}}
}{
r_{ij}^{\mathrm{sub}}
}.
\end{equation}

The interface stiffness is
\begin{equation}
c_{\mathrm{int}}
=
\alpha_{\mathrm{int}}c,
\end{equation}
with
\begin{equation}
\alpha_{\mathrm{int}} = 1.
\end{equation}
Only virtual substrate points satisfying
\begin{equation}
r_{ij}^{0,\mathrm{sub}} \le \delta
\end{equation}
are included.

The interface force density contribution is
\begin{equation}
p_i^{\mathrm{sub}}
=
\sum_j
c_{\mathrm{int}}
s_{ij}^{\mathrm{sub}}
n_{ij}^{\mathrm{sub}}
V_j.
\end{equation}
The total force density is
\begin{equation}
p_i
=
p_i^{\mathrm{bar}}
+
p_i^{\mathrm{sub}}.
\end{equation}

\section{Boundary Conditions}

The simulation uses affine displacement boundary conditions. At outer load step \(k\), the imposed macroscopic strain is
\begin{equation}
\varepsilon_k = k\Delta\varepsilon.
\end{equation}
The strain increment is
\begin{equation}
\Delta\varepsilon
=
\frac{\varepsilon_{\mathrm{target}}}{N_{\mathrm{load}}},
\end{equation}
where currently
\begin{equation}
\varepsilon_{\mathrm{target}} = 0.004,
\qquad
N_{\mathrm{load}} = 80.
\end{equation}

The first and last three nodes are displacement-prescribed:
\begin{equation}
u_i = \varepsilon_k x_i,
\qquad
i=0,1,2,
\end{equation}
and
\begin{equation}
u_i = \varepsilon_k x_i,
\qquad
i=N-3,N-2,N-1.
\end{equation}
Thus the left boundary satisfies
\begin{equation}
u(0)=0,
\end{equation}
but the neighboring grip nodes follow the same affine field:
\begin{equation}
u(x)=\varepsilon_k x.
\end{equation}
The right boundary satisfies
\begin{equation}
u(L)=\varepsilon_k L.
\end{equation}

The velocities of the prescribed nodes are set to zero:
\begin{equation}
v_i = 0.
\end{equation}
The imposed displacement is re-applied during the inner relaxation loop:
\begin{equation}
u_i = \varepsilon_k x_i.
\end{equation}

\section{Outer Load Iterations}

The simulation is quasi-static and is advanced through outer load increments. At load step \(k\),
\begin{equation}
\varepsilon_k
=
\varepsilon_{k-1}
+
\Delta\varepsilon.
\end{equation}

The converged displacement field from the previous load step is retained:
\begin{equation}
u_i^{k,0}
=
u_i^{k-1,\mathrm{conv}}
\end{equation}
for all unconstrained interior nodes.

Only the prescribed boundary nodes are updated to the new affine displacement:
\begin{equation}
u_i^{k,0}
=
\varepsilon_k x_i.
\end{equation}
The artificial dynamic relaxation state is reset:
\begin{equation}
v_i^{k,0}=0.
\end{equation}
Thus the solution path is continuous in displacement, while each equilibrium solve starts from zero velocity.

\section{Inner Dynamic Relaxation}

For each outer load increment, an inner dynamic relaxation loop solves the equilibrium problem. The relaxation equation is
\begin{equation}
\rho \ddot{u}_i
+
c_d \dot{u}_i
=
p_i+b_i,
\end{equation}
where \(b_i\) is the external body force density. In the current simulations,
\begin{equation}
b_i=0.
\end{equation}
The density and damping coefficient are
\begin{equation}
\rho = 1,
\qquad
c_d = 5.
\end{equation}

The acceleration is
\begin{equation}
a_i^n
=
\frac{
p_i^n+b_i-c_d v_i^n
}{
m_i
}.
\end{equation}
The mass-like coefficient is
\begin{equation}
m_i = \rho.
\end{equation}
The velocity update is explicit:
\begin{equation}
v_i^{n+1}
=
v_i^n+a_i^n\Delta t.
\end{equation}
The displacement update is
\begin{equation}
u_i^{n+1}
=
u_i^n+v_i^{n+1}\Delta t.
\end{equation}
The time step is
\begin{equation}
\Delta t = 5\times 10^{-4}.
\end{equation}

For prescribed boundary nodes,
\begin{equation}
v_i^{n+1}=0,
\end{equation}
and
\begin{equation}
u_i^{n+1}=\varepsilon_k x_i.
\end{equation}

The convergence criterion is based on the maximum velocity:
\begin{equation}
\max_i |v_i| \le \mathrm{tol},
\end{equation}
with
\begin{equation}
\mathrm{tol}=10^{-6}.
\end{equation}

\section{Linear Softening Damage}

Damage is active with the linear softening model:
\begin{equation}
\texttt{degradation\_model}=1.
\end{equation}

The critical elastic stretch is
\begin{equation}
s_0
=
\sqrt{
\frac{3G_c}{E\delta}
}.
\end{equation}
The fracture energy parameter is
\begin{equation}
G_c = 10^{-8}.
\end{equation}
The softening parameter is
\begin{equation}
R_d = G_c.
\end{equation}
The complete-failure stretch is
\begin{equation}
s_c
=
1.5
\left(
s_0
+
\frac{2R_d}{cs_0}
\right).
\end{equation}

For each bond, the historical maximum stretch is
\begin{equation}
s_{ij}^{\max}
=
\max_{\tau\le t}s_{ij}(\tau).
\end{equation}
If
\begin{equation}
s_{ij}^{\max} \le s_0,
\end{equation}
the bond is elastic and
\begin{equation}
\mu_{ij}=1.
\end{equation}
If
\begin{equation}
s_0 < s_{ij}^{\max} < s_c,
\end{equation}
the bond is softened with
\begin{equation}
\mu_{ij}
=
\frac{s_c-s_{ij}^{\max}}{s_c-s_0}.
\end{equation}
The envelope force magnitude is
\begin{equation}
f_{\mathrm{env}}
=
cs_0\mu_{ij}.
\end{equation}

During continued loading on the envelope, the bond force is
\begin{equation}
f_{ij}
=
f_{\mathrm{env}} n_{ij} w_{ij}.
\end{equation}
During unloading or reloading, the degraded secant stiffness is
\begin{equation}
k_{\mathrm{unload}}
=
\frac{f_{\mathrm{env}}}{s_{ij}^{\max}},
\end{equation}
and the bond force is
\begin{equation}
f_{ij}
=
k_{\mathrm{unload}}s_{ij}n_{ij}w_{ij}.
\end{equation}
If
\begin{equation}
s_{ij}^{\max}\ge s_c,
\end{equation}
the bond is fully failed:
\begin{equation}
f_{ij}=0.
\end{equation}

\section{Damage Zone}

Damage is allowed only in the central active region of the bar:
\begin{equation}
|x_i-0.5| < 0.45.
\end{equation}
Since the bar occupies \(x\in[0,1]\), this corresponds approximately to
\begin{equation}
0.05 < x_i < 0.95.
\end{equation}
This prevents damage directly inside the displacement-controlled grip regions.

\section{Nodal Damage Variable}

The scalar nodal damage variable is computed from the remaining bond integrity around node \(i\):
\begin{equation}
D_i
=
1
-
\frac{
\sum_{j\in\mathcal{H}_i}\mu_{ij}V_j
}{
\sum_{j\in\mathcal{H}_i}V_j
}.
\end{equation}
Thus
\begin{equation}
D_i=0
\end{equation}
means all connected bonds are intact, while
\begin{equation}
D_i=1
\end{equation}
means all connected bonds are fully failed.

\section{Strain Calculation}

The output strain is a post-processed local nodal strain, not the peridynamic bond stretch. The element strain is
\begin{equation}
\varepsilon_{i+1/2}
=
\frac{u_{i+1}-u_i}{x_{i+1}-x_i}.
\end{equation}
For a uniform grid,
\begin{equation}
\varepsilon_{i+1/2}
=
\frac{u_{i+1}-u_i}{\Delta x}.
\end{equation}
The nodal strain is
\begin{equation}
\varepsilon_0
=
\varepsilon_{1/2},
\end{equation}
\begin{equation}
\varepsilon_{N-1}
=
\varepsilon_{N-3/2},
\end{equation}
and, for interior nodes,
\begin{equation}
\varepsilon_i
=
\frac{1}{2}
\left(
\varepsilon_{i-1/2}
+
\varepsilon_{i+1/2}
\right).
\end{equation}
Equivalently,
\begin{equation}
\varepsilon_i
=
\frac{1}{2}
\left[
\frac{u_i-u_{i-1}}{\Delta x}
+
\frac{u_{i+1}-u_i}{\Delta x}
\right].
\end{equation}

\section{Stress Calculation}

The stress-like quantity is computed from peridynamic bond force moments:
\begin{equation}
\sigma_i
=
C_\sigma
\sum_{j\in\mathcal{H}_i}
f_{ij}\xi_{ij}.
\end{equation}
The coefficient used is
\begin{equation}
C_\sigma=\frac{1}{2}.
\end{equation}
This is a peridynamic virial-type stress measure. Before damage, a homogeneous affine solution should satisfy approximately
\begin{equation}
\varepsilon_i \approx \varepsilon_{\mathrm{applied}}.
\end{equation}
After damage begins, the stress is no longer simply \(E\varepsilon\), because the bond stiffness is degraded through \(\mu_{ij}\).

\section{Energy Quantities}

The elastic strain energy contribution of an intact bond is
\begin{equation}
W_{ij}
=
\frac{1}{4}
c s_{ij}^2 |\xi_{ij}| w_{ij}.
\end{equation}
The factor \(1/4\) appears because pairwise bond energy is shared between interacting material points.

For softened unloading or reloading,
\begin{equation}
W_{ij}
=
\frac{1}{4}
k_{\mathrm{unload}}
s_{ij}^2
|\xi_{ij}|
w_{ij}.
\end{equation}

The substrate energy is
\begin{equation}
W_i^{\mathrm{sub}}
=
\sum_j
\frac{1}{4}
c_{\mathrm{int}}
\left(s_{ij}^{\mathrm{sub}}\right)^2
r_{ij}^{0,\mathrm{sub}}
V_j.
\end{equation}

The kinetic energy is
\begin{equation}
K_i
=
\frac{1}{2}m_i v_i^2.
\end{equation}
The damping dissipation is accumulated as
\begin{equation}
\mathcal{D}_i^{n+1}
=
\mathcal{D}_i^n
+
m_i c_d v_i^2\Delta t.
\end{equation}
The total nodal energy density is stored as
\begin{equation}
E_i^{\mathrm{tot}}
=
W_i
+
W_i^{\mathrm{sub}}
+
K_i
+
\mathcal{D}_i
+
G_i^{\mathrm{fracture}}.
\end{equation}

\section{Output Fields}

For every outer load increment, the code writes
\begin{equation}
\texttt{result/conf}k\texttt{.dat}.
\end{equation}
Each row contains
\begin{equation}
x_i,
\qquad
u_i,
\qquad
\varepsilon_i.
\end{equation}
OVITO files are written separately into
\begin{equation}
\texttt{ovito\_files/}.
\end{equation}
They contain position, displacement, damage, stress, and energy-related fields for visualization.

\section{Algorithm Summary}

At initialization, the grid is constructed as
\begin{equation}
x_i=i\Delta x.
\end{equation}
The families are built using
\begin{equation}
|x_j-x_i|\le\delta+10^{-12}.
\end{equation}
Surface correction factors \(G_i\) are computed, and the grip-node corrections are set to unity.

For every outer load step \(k\):
\begin{align}
\varepsilon_k &= \varepsilon_{k-1}+\Delta\varepsilon, \\
u_i^{k,0} &= u_i^{k-1,\mathrm{conv}}
\quad \text{for interior nodes}, \\
u_i^{k,0} &= \varepsilon_k x_i
\quad \text{for prescribed boundary nodes}, \\
v_i^{k,0} &= 0.
\end{align}

For every inner dynamic relaxation iteration:
\begin{enumerate}
\item compute bar-bar bond stretches \(s_{ij}\);
\item update the historical maximum stretches \(s_{ij}^{\max}\);
\item apply the elastic, softened, or failed bond law;
\item compute substrate/interface forces;
\item sum the total force density \(p_i\);
\item update velocity and displacement;
\item reimpose affine boundary displacements;
\item compute strain, stress, damage, and energies;
\item stop when
\begin{equation}
\max_i |v_i| \le 10^{-6}.
\end{equation}
\end{enumerate}

After convergence, the result files are written and the next outer load increment begins.

\end{document}
